% !TeX root = ../manuscript.tex


\chapter{图表示例}
\section{测试}
This is a test for compiling.

It is slower than the \LaTeX mk, a lot of slow

what about the `neovim`, it is faster than the texifier.





\section{插图}

图片通常在 \env{figure} 环境中使用 \cs{includegraphics} 插入，如图~\ref{fig:example} 的源代码。
建议矢量图片使用 PDF 格式，比如数据可视化的绘图；
照片应使用 JPG 格式；
其他的栅格图应使用无损的 PNG 格式。
注意，LaTeX 不支持 TIFF 格式；EPS 格式已经过时。

\begin{figure}
  \centering
  \includegraphics[width=0.5\linewidth]{example-image-a.pdf}
  \caption*{国外的期刊习惯将图表的标题和说明文字写成一段，需要改写为标题只含图表的名称，其他说明文字以注释方式写在图表下方，或者写在正文中。}
  \caption{示例图片标题}
  \label{fig:example}
\end{figure}

图应有编号。由“图”和从“1”开始的阿拉伯数字组成，例如“图1”等，图较多时，可分章编号。
图应有图题，图题即图的名称，置于图的编号之后。图的编号和图题应置于图下方。
若图或表中有附注，采用英文小写字母顺序编号，附注写在图或表的下方。
国外的期刊习惯将图表的标题和说明文字写成一段，需要改写为标题只含图表的名称，其他说明文字以注释方式写在图表下方，或者写在正文中。

\begin{figure}
  \centering
  \subcaptionbox{分图 A\label{fig:subfig-a}}
    {\includegraphics[width=0.35\linewidth]{example-image-a.pdf}}
  \subcaptionbox{分图 B\label{fig:subfig-b}}
    {\includegraphics[width=0.35\linewidth]{example-image-b.pdf}}
  \caption{多个分图的示例}
  \label{fig:multi-image}
\end{figure}

如果一个图由两个或两个以上分图组成时，各分图分别以 (a)、(b)、(c)...... 作为图序，并须有分图题。
推荐使用 \pkg{subcaption} 宏包来处理， 比如图~\ref{fig:subfig-a} 和图~\ref{fig:subfig-b}。

\section{表格}

表应具有自明性。为使表格简洁易读，尽可能采用三线表，如表~\ref{tab:three-line}。
三条线可以使用 \pkg{booktabs} 宏包提供的命令生成。

\begin{table}
  \centering
  \caption{三线表示例}
  \label{tab:three-line}
  \begin{tabular}{ll}
    \toprule
    文件名          & 描述                         \\
    \midrule
    bnuthesis.dtx   & 模板的源文件，包括文档和注释 \\
    bnuthesis.cls   & 模板文件                     \\
    bnuthesis-*参考文献.bst & BibTeX 参考文献表样式文件    \\
    \bottomrule
  \end{tabular}
\end{table}


表应有编号。由“表”和从“1”开始的阿拉伯数字组成，例如“表1”等，表较多时，可分章编号。
表应有表题，表题即表的名称，置于表的编号之后。表的编号和表题应置于表上方。
表格如果有附注，尤其是需要在表格中进行标注时，可以使用 \pkg{threeparttable} 宏包。

\begin{table}
  \centering
  \begin{threeparttable}[c]
    \caption{带附注的表格示例}
    \label{tab:three-part-table}
    \begin{tabular}{ll}
      \toprule
      文件名                 & 描述                         \\
      \midrule
      bnuthesis.dtx\tnote{a} & 模板的源文件，包括文档和注释 \\
      bnuthesis.cls\tnote{b} & 模板文件                     \\
      bnuthesis-*.bst        & BibTeX 参考文献表样式文件    \\
      \bottomrule
    \end{tabular}
    \begin{tablenotes}
      \item [a] 可以通过 xelatex 编译生成模板的使用说明文档；
        使用 xetex 编译 \file{bnuthesis.ins} 时则会从 \file{.dtx} 中去除掉文档和注释，得到精简的 \file{.cls} 文件。
      \item [b] 更新模板时，一定要记得编译生成 \file{.cls} 文件，否则编译论文时载入的依然是旧版的模板。
    \end{tablenotes}
  \end{threeparttable}
\end{table}

如某个表需要转页接排，可以使用 \pkg{longtable} 宏包，需要在随后的各页上重复表的编号。
编号后跟表题（可省略）和“（续）”，置于表上方。续表均应重复表头。

\begin{longtable}{cccc}
    \caption{跨页长表格的表题}
    \label{tab:longtable} \\
    \toprule
    表头 1 & 表头 2 & 表头 3 & 表头 4 \\
    \midrule
  \endfirsthead
    \caption*{续表~\thetable\quad 跨页长表格的表题} \\
    \toprule
    表头 1 & 表头 2 & 表头 3 & 表头 4 \\
    \midrule
  \endhead
    \bottomrule
  \endfoot
\end{longtable}


\section{公式}

\subsection{勾股定理}
勾股定理是欧几里得几何中的一个基本定理，描述了直角三角形三边之间的关系。它表明，在一个直角三角形中，直角边的平方和等于斜边的平方。这可以用公式表示为：
\begin{equation}\label{eq:pythagorean}
    a^2 + b^2 = c^2
\end{equation}

其中，\(a\) 和 \(b\) 是直角三角形的两个直角边，而 \(c\) 是斜边。
根据公式~\eqref{eq:pythagorean}，当直角边的长度已知时我们可以计算出斜边的长度。

\subsection{格林公式}
格林公式是一个在向量分析中非常重要的定理，它提供了一个将双重积分转换为曲线积分的方法。具体来说，格林公式将平面区域上的双重积分转换为该区域边界上的曲线积分。

格林公式的数学表达为：
\begin{equation}\label{eq:greens_theorem}
    \oint_C (L \dif x + M \dif y) = \iint_D \left(\frac{\partial M}{\partial x} - \frac{\partial L}{\partial y}\right) \dif A
\end{equation}
其中，\(C\) 是平面上的简单闭曲线，\(D\) 是由 \(C\) 围成的区域，\(L\) 和 \(M\) 是定义在 \(D\) 上的连续可微函数，\(\dif x\)和 \(\dif y\) 分别是曲线 \(C\) 上点的横纵坐标的微分。



\section{算法}

算法环境可以使用 \pkg{algorithms} 或者 \pkg{algorithm2e} 宏包。

\renewcommand{\algorithmicrequire}{\textbf{输入：}\unskip}
\renewcommand{\algorithmicensure}{\textbf{输出：}\unskip}

\begin{algorithm}
  \caption{Calculate $y = x^n$}
  \label{alg1}
  \small
  \begin{algorithmic}
    \REQUIRE $n \geq 0$
    \ENSURE $y = x^n$

    \STATE $y \leftarrow 1$
    \STATE $X \leftarrow x$
    \STATE $N \leftarrow n$

    \WHILE{$N \neq 0$}
      \IF{$N$ is even}
        \STATE $X \leftarrow X \times X$
        \STATE $N \leftarrow N / 2$
      \ELSE[$N$ is odd]
        \STATE $y \leftarrow y \times X$
        \STATE $N \leftarrow N - 1$
      \ENDIF
    \ENDWHILE
  \end{algorithmic}
\end{algorithm}



\section{题目\texorpdfstring{\(e=mc^2\)}{e=mc2}}

当章节出现数学公式时，请使用上述格式。

\begin{theorem}[费马大定理]
对于任何大于2的正整数 $n$，不存在正整数 $x, y, z$ 使得 $x^n + y^n = z^n$ 成立。
\end{theorem}

\begin{proof}
假设存在这样的 $x, y, z$，使得 $x^n + y^n = z^n$ 成立。不失一般性，假设 $x, y, z$ 互素。

考虑模 $p$ 的情况，其中 $p$ 是任意的质数。

\begin{enumerate}
\item 当 $n = 2$ 时，费马大定理是勾股定理，已经被证明。所以我们假设 $n > 2$。
\item 当 $n$ 是偶数时，$x^n + y^n$ 和 $z^n$ 都可以被因式分解为两个平方数的和。但这违背了费马的最后定理的要求，因此不存在这样的解。
\item 当 $n$ 是奇数时，我们考虑同样的方式，由于 $x^n + y^n$ 是一个奇数，而 $z^n$ 是一个偶数，所以这也不可能成立。
\end{enumerate}

因此，我们得出结论，费马大定理成立。
\end{proof}
